% optimization/alignment/30_sqp_concrete.tex

\section{Sequential Quadratic Programming (SQP)}

\subsection{Problem Structure}

We consider the discrete optimization problem
\[
\min_{\mathbf{x}} J(\mathbf{x})
\quad \text{subject to} \quad
g(\mathbf{x}) = 0,\;
h(\mathbf{x}) \le 0.
\]

\begin{definition}[Parameter vector]
\[
\mathbf{x} =
(\kappa_0,\dots,\kappa_N,\phi_0,\dots,\phi_N).
\]
\end{definition}

% ============================================================
\subsection{Residual Formulation}
% ============================================================

\begin{definition}[Residual vector]
\[
J(\mathbf{x}) = \frac{1}{2} \| \mathbf{r}(\mathbf{x}) \|^2.
\]
\end{definition}

Typical residual components:
\[
r_i^{(\kappa)} = \sqrt{\gamma}
\frac{\kappa_{i+1}-\kappa_i}{\Delta s},
\quad
r_i^{(\phi)} = \sqrt{\delta}
\frac{\phi_{i+1}-\phi_i}{\Delta s},
\]
\[
r_i^{(\mathrm{dyn})}
=
\sqrt{\beta}
(v^2 \kappa_i - g \sin\phi_i).
\]

% ============================================================
\subsection{Gauss--Newton Step}
% ============================================================

\begin{definition}[SQP step]
At iteration \(k\), compute \(\mathbf{p}_k\) from:
\[
(J_r^T J_r)\mathbf{p}_k = -J_r^T \mathbf{r}.
\]
\end{definition}

The Hessian is approximated by:
\[
H \approx J_r^T J_r.
\]

% ============================================================
\subsection{Sparse Structure}
% ============================================================

The system matrix has band structure:
\begin{itemize}
\item tridiagonal coupling for smoothness terms,
\item diagonal contributions from dynamic terms.
\end{itemize}

% ============================================================
\subsection{Constraints Handling}
% ============================================================

Constraints are handled via:
\begin{itemize}
\item projection,
\item penalty terms,
\item Lagrange multipliers.
\end{itemize}

% ============================================================
\subsection{Iteration Scheme}
% ============================================================

\begin{algorithm}
\caption{SQP iteration}
\begin{algorithmic}
\State initialize $\mathbf{x}_0$
\For{$k = 0,1,\dots$}
  \State compute $\mathbf{r}(\mathbf{x}_k)$
  \State compute $J_r(\mathbf{x}_k)$
  \State solve $(J_r^T J_r)\mathbf{p}_k = -J_r^T \mathbf{r}$
  \State choose step size $\alpha_k$
  \State update $\mathbf{x}_{k+1} = \mathbf{x}_k + \alpha_k \mathbf{p}_k$
\EndFor
\end{algorithmic}
\end{algorithm}

% ============================================================
\subsection{Interpretation}
% ============================================================

\begin{summary}
SQP provides a practical and efficient method for solving the discrete
alignment optimization problem.

Its efficiency follows from the sparse and local structure of the
underlying model.
\end{summary}
